\section{Análisis de los datos de salida}

\subsection{Resultados por escenario}

\begin{table}[H]
\centering
\small
\begin{tabular}{lcccc}
\toprule
 & \multicolumn{2}{c}{\textbf{Escenario base}} & \multicolumn{2}{c}{\textbf{Escenario alternativo}} \\
\cmidrule(lr){2-3}\cmidrule(lr){4-5}
\textbf{Indicador} & \textbf{Media} & \textbf{Desvío} & \textbf{Media} & \textbf{Desvío} \\
\midrule
Espera hasta el triage      & --- & --- & --- & --- \\
Espera hasta el consultorio & --- & --- & --- & --- \\
Espera de pacientes verdes  & --- & --- & --- & --- \\
Espera de urgencias reales  & --- & --- & --- & --- \\
Permanencia total           & --- & --- & --- & --- \\
Pacientes en sala de espera & --- & --- & --- & --- \\
\bottomrule
\end{tabular}
\caption{Resultados sobre 30 réplicas por escenario.}
\label{tab:resultados}
\end{table}

\subsection{Test de medias}

Para determinar si las diferencias observadas entre escenarios resultan
estadísticamente significativas se aplica un test de diferencia de medias
para muestras independientes. Las réplicas de cada escenario se generan con
semillas distintas e independientes entre sí, por lo que las dos muestras son
independientes. No se asume igualdad de varianzas, de modo que se emplea el
test de Welch.

Sean $\bar{x}_1$ y $\bar{x}_2$ las medias muestrales de cada escenario,
$s_1^2$ y $s_2^2$ sus varianzas muestrales y $n_1 = n_2 = 30$ el número de
réplicas. Las hipótesis son

\[
H_0: \mu_1 = \mu_2
\qquad\text{frente a}\qquad
H_1: \mu_1 \neq \mu_2
\]

con un nivel de significancia $\alpha = 0{,}05$.

El estadístico de contraste es

\begin{equation}
t = \frac{\bar{x}_1 - \bar{x}_2}
         {\sqrt{\dfrac{s_1^2}{n_1} + \dfrac{s_2^2}{n_2}}}
\label{eq:t}
\end{equation}

y los grados de libertad se aproximan mediante la ecuación de
Welch--Satterthwaite:

\begin{equation}
\nu = \frac{\left(\dfrac{s_1^2}{n_1} + \dfrac{s_2^2}{n_2}\right)^{\!2}}
           {\dfrac{\left(s_1^2/n_1\right)^2}{n_1 - 1} +
            \dfrac{\left(s_2^2/n_2\right)^2}{n_2 - 1}}
\label{eq:welch}
\end{equation}

Se rechaza $H_0$ cuando $|t| > t_{\nu,\,1-\alpha/2}$.

El intervalo de confianza del $100(1-\alpha)\%$ para la diferencia de medias
resulta

\begin{equation}
(\bar{x}_1 - \bar{x}_2) \;\pm\;
t_{\nu,\,1-\alpha/2}\,
\sqrt{\frac{s_1^2}{n_1} + \frac{s_2^2}{n_2}}
\label{eq:ic}
\end{equation}

El intervalo de confianza resulta más informativo que el valor $p$ para los
fines de este trabajo, por cuanto expresa la magnitud de la mejora en las
mismas unidades del indicador y permite juzgar su relevancia operativa, y no
solamente su significancia estadística.

\begin{table}[H]
\centering
\small
\begin{tabular}{lccccc}
\toprule
\textbf{Indicador} & \textbf{Diferencia} & $t$ & $\nu$ & \textbf{IC 95\%} & \textbf{Decisión} \\
\midrule
Espera hasta el triage      & --- & --- & --- & --- & --- \\
Espera hasta el consultorio & --- & --- & --- & --- & --- \\
Espera de pacientes verdes  & --- & --- & --- & --- & --- \\
Espera de urgencias reales  & --- & --- & --- & --- & --- \\
Permanencia total           & --- & --- & --- & --- & --- \\
Pacientes en sala de espera & --- & --- & --- & --- & --- \\
\bottomrule
\end{tabular}
\caption{Test de medias entre escenarios. Nivel de significancia 0,05.}
\label{tab:test}
\end{table}

\begin{pendiente}
Completar las tablas~\ref{tab:resultados} y~\ref{tab:test} con los resultados
de las corridas de producción. Las fórmulas~\eqref{eq:t},
\eqref{eq:welch} y~\eqref{eq:ic} se aplican sobre las sesenta observaciones
exportadas desde AnyLogic.
\end{pendiente}

\subsection{Interpretación}

El diseño del estudio permite distinguir tres desenlaces posibles, cada uno
con una recomendación distinta:

\begin{table}[H]
\centering
\small
\begin{tabular}{p{3.6cm}p{5.0cm}p{5.0cm}}
\toprule
\textbf{Desenlace} & \textbf{Evidencia que lo identifica} & \textbf{Implicancia} \\
\midrule
El cuello de botella se resuelve &
Cae la espera de triage y la permanencia total disminuye en magnitud
comparable. La espera de consultorio aumenta poco. &
El triage era la restricción efectiva. La incorporación se justifica por sí
sola. \\
El cuello de botella se traslada &
Cae la espera de triage pero la de consultorio aumenta en magnitud similar.
La permanencia total apenas se modifica. &
La restricción son los consultorios. La medida mejora la oportunidad de la
clasificación, no la congestión. \\
El sistema se satura aguas abajo &
La cola de consultorios crece de manera sostenida durante toda la ventana de
medición. &
La capacidad médica es insuficiente para el caudal. Corresponde evaluar un
quinto consultorio. \\
\bottomrule
\end{tabular}
\caption{Desenlaces posibles del estudio y su interpretación.}
\end{table}

Cabe destacar que incluso en el segundo desenlace la intervención conserva
valor clínico, puesto que desplaza la espera desde una cola ciega hacia una
cola priorizada: el paciente grave espera menos aunque el promedio del
conjunto se modifique poco. La ausencia de diferencia significativa en el
tiempo total constituye, en ese caso, un resultado válido y relevante.

